\section{The CIPHER Benchmark}
\label{sec:benchmark}

\subsection{Difficulty levels}

\begin{table}[h]
\centering\small
\caption{CIPHER difficulty levels (50 samples each, $N{=}250$ total).}
\label{tab:levels}
\begin{tabular}{lcccp{5cm}}
\toprule
Level & Key vis. & Ops & NL ambiguity & Primary challenge \\
\midrule
L1 & High   & 1 & Exact         & Single-op baseline \\
L2 & Medium & 2 & Exact         & Two-op chaining \\
L3 & Low    & 2 & Paraphrase    & Surface-form variation \\
L4 & Low    & 3 & Reversed      & Reversed description \\
L5 & Min.   & 4 & Compositional & Underspecified instruction \\
\bottomrule
\end{tabular}
\end{table}

\paragraph{Key visibility} varies font size (12\%$\to$5\% of frame height),
alpha (255$\to$140), and exposure duration (4.5$\to$2.5s).

\paragraph{NL ambiguity templates:}
(1)~exact canonical; (2)~paraphrase from pool of 5 per op;
(3)~reversed order with ``undo'' framing;
(4)~compositional: ``The video has been edited. Undo all transformations.''

\paragraph{Operation Vocabulary.}
The operational vocabulary consists of seven foundational transformations: \texttt{reverse}, \texttt{hflip}, \texttt{vflip}, \texttt{rotate90}, \texttt{rotate270}, \texttt{speed2x}, and \texttt{grayscale}. This specific set was chosen because these operations are visually distinct, non-destructive (except \texttt{grayscale}, used primarily as a distractor), and strictly composable via standard \texttt{ffmpeg} filters. They demand different types of reasoning: spatial (\texttt{hflip}, \texttt{rotate}), temporal (\texttt{reverse}, \texttt{speed2x}), and color-based (\texttt{grayscale}), providing a comprehensive testbed for multimodal execution without requiring complex parameter tuning from the model.

\subsection{Procedural generation and Dataset Integrity}

Samples are fully deterministic from $(seed, level)$.
Background clips are drawn uniformly from a Kinetics-400 pool
(diverse human activities, minimal on-screen text).
Test split: seeds 0--49. Challenge split: seeds 50--99.

A persistent challenge in evaluating large language and vision models is the risk of data contamination, where static benchmark data inadvertently leaks into the pre-training corpus. CIPHER mitigates this risk entirely through programmatic dataset synthesis. The benchmark is procedurally generated from a tuple of $(seed, \text{level})$, which deterministically selects the background clip, the randomly generated 7-character key, the sequence of operations, and the natural language instruction template. Because samples are dynamically constructed on-demand using disjoint background video pools, researchers can regenerate the benchmark with novel seeds, guaranteeing an uncontaminated evaluation environment and an effectively infinite supply of test instances.

\subsection{Metrics}
\label{sec:metrics}

\paragraph{EM.} $\frac{1}{N}\sum_i \mathbf{1}[\hat{k}_i = k^*_i]$.

\paragraph{CharF1.} Character-level precision/recall F1 (partial credit).

\paragraph{PipeF1.}
\begin{equation}
\text{PipeF1} = 0.25\cdot\overline{\text{S1}} + 0.25\cdot\overline{\text{S2}} + 0.50\cdot\overline{\text{CharF1}}
\end{equation}
S3 receives double weight as the terminal objective;
S1/S2 serve as diagnostic leading indicators.